\section{Related work}\label{sec:related_work}

\emph{Dispersion problems.} \citet{rosenkrantz2000} introduced dispersion under capacity and cost constraints. The CDP was brought to the maximum-diversity benchmark setting by \citet{peiro2021}, who proposed a GRASP--VND hybrid within strategic oscillation and built the 100-instance testbed from MDPLIB (capacities in [1, 1000], B equal to 0.2 or 0.3 of the total capacity), which is the testbed we use. \citet{marti2021b} added a node-packing-based model and a scatter-search heuristic; \citet{landete2023} studied further formulations and valid inequalities; \citet{mladenovic2022b} proposed basic, general and skewed variable neighbourhood search; \citet{lu2023} a solution-based tabu search (SBTS); and \citet{rosati2024} a multi-neighbourhood simulated annealing together with a harder dataset from real maps and compact models. \citet{gomez2022} proposed a multi-start biased-randomised algorithm, and later work added stochastic and learning-based variants \citep{gomez2023,ghorbani2024}. The GDP was introduced by \citet{martinezgavara2021} with a GRASP and a tabu search; \citet{lozanoosorio2022} gave a compact and a standard linear formulation and a multi-start basic VNS, tightened the budgets of the benchmark (whose cost constraint was inactive at the optimum) and released instances, code and solutions; \citet{mladenovic2022a} proposed a multiple-neighbourhood tabu search.

\emph{Multi-objective diversity.} Multi-objective diversity has been studied since 2018. \citet{vera2018} proposed a MOEA for the simultaneous optimisation of up to five diversity measures (Max-Sum, Max-Min, Max-MinSum, Min-Diff, Min-P-center), with fixed and free subset size; \citet{colmenar2018} defined the bi-objective diversity problem (Max-Sum and Max-Min), compared NSGA-II, SPEA2, GRASP, iterated greedy, tabu search and VNS on the 145 GKD instances of MDPLIB (75 small, 50 medium and 20 large, with n up to 500 and m up to 50), each method tuned with irace on 29 of the instances (at most 500 experiments per algorithm; more experiments gave similar results) and judged by hypervolume (computed against the origin), set coverage and the $\varepsilon$-indicator, following the guidelines of \citet{knowles2006}, with 10 runs per instance. The tabu search was the best method overall: it reached the best hypervolume (ties included) on 109 of the 145 instances (17 of the 20 large ones), the best average $\varepsilon$-indicator and the shortest run times, while GRASP had the highest average hypervolume over all instances (and the best on the medium ones) and NSGA-II the highest set coverage; the tabu search owed its hypervolume to the extreme regions of the front, whereas NSGA-II concentrated on the central part. They are the direct precedent of our MDP experiments. Their study uses the GKD family only, no decomposition, adaptive operator selection or objective scaling, and no constrained variant. The tutorial of \citet{marti2023} has a chapter on multi-objective diversity. Exact approaches exist for related bi-objective location--diversity problems: a bi-objective p-median/max-sum diversity problem \citep{yang2026} and a bi-criteria p-hub median and dispersion problem \citep{ramamoorthy2024}. We found no prior work that evaluates evolutionary multi-objective algorithms on the MaxSum + MaxMin versions of the \emph{capacitated} and \emph{generalised} dispersion problems, nor one that studies operator-pool design and adaptive operator selection in the diversity setting; the search covered web and publisher indices on 5 October 2026 \todo{repeat the search (Scopus/Web of Science) just before submitting}.

\emph{Adaptive operator selection in MOEAs.} Bandit-based AOS was analysed by \citet{fialho2010}, who proposed a sliding-window UCB and studied it on artificial reward scenarios and on a genetic algorithm for the Royal Road problem, including its sensitivity to the hyper-parameters, and by \citet{dacosta2008} (dynamic bandits with a change-point test). In MOEA/D, \citet{li2014} introduced FRRMAB over differential-evolution operators and reported, on the UF benchmark, that it improves on other AOS strategies and on uniformly random operator selection, which was the worst of the strategies compared; \citet{sun2020} replaced the bandit by a dynamic Thompson sampler (19 benchmark problems), and \citet{geng2023} by a classification tree that predicts the outcome of each operator (UF and LZ problems). These studies use continuous test problems and pools of differential-evolution variants. Our setting is a discrete one with a pool of complementary neighbourhood operators; there, uniform selection from the pool is as good as any selector we tried on instances up to about 200 elements, and the circumstances in which selectors win are of a different kind (Section 6.11).
